% =====================================================================
%  6. Salvaguardas ambientais e sociais
% =====================================================================
\section{Salvaguardas ambientais e sociais}
\label{sec:salvaguardas}

Todas as fases descritas nas seções anteriores observarão o Quadro Ambiental
e Social do Banco Mundial \cite{bm-esf}, aplicado ao PSH-PR por meio de cinco
instrumentos. O Marco de Gestão Ambiental e Social (MGAS) orienta o controle
dos impactos das obras e da operação. O Plano de Engajamento das Partes
Interessadas (PEPI) orienta a participação das comunidades. Os Procedimentos
de Gestão de Mão de Obra (PGR) tratam da saúde e da segurança dos
trabalhadores. O Marco de Política de Reassentamento (MRF) rege as questões
de posse e uso da terra, e o Mecanismo de Queixas e Reclamações (MQR) recebe e
responde às manifestações da população.

Os riscos variam conforme a fase. No planejamento, o principal risco é
excluir grupos vulneráveis ou comunidades tradicionais dos diagnósticos, o
que será mitigado com consultas públicas e metodologias participativas na
seleção das áreas, além da busca ativa de famílias em áreas remotas e de
reuniões comunitárias em linguagem simples. Na captação e na distribuição de
água, os riscos são a contaminação de aquíferos e os impactos temporários das
obras, como ruído, poeira e erosão, controlados pela adoção de normas
técnicas de perfuração, pelo gerenciamento de resíduos e pela umectação das
vias. Como o poço e o reservatório ocupam áreas específicas, a titularidade
dessas áreas será verificada previamente e formalizada por anuências e termos
de servidão, conforme o MRF. Na operação, os riscos são a falta de
sustentabilidade técnica e os conflitos pelo uso da água, que o modelo de
gestão, a tarifa social e o treinamento das associações procuram evitar. No
esgotamento, os riscos são o manejo inadequado do lodo e a contaminação do
solo por instalação incorreta, enfrentados com testes de estanqueidade e
capacitação dos proprietários. Nas obras em geral, protocolos de saúde e
segurança do trabalho protegerão operários, técnicos e famílias. O MQR será
implantado e divulgado nas comunidades, e seus tempos de resposta serão
monitorados semestralmente \cite[Quadros 25 e 26]{mop2026}. O
\autoref{qua:salvaguardas} organiza esses riscos por atividade.

\begin{quadro}[htbp]
\caption{Riscos ambientais e sociais e medidas de mitigação do Subcomponente 3.2.b}
\label{qua:salvaguardas}
\footnotesize
\renewcommand{\arraystretch}{1.25}
\begin{tabularx}{\textwidth}{|P{2.4cm}|L|L|P{2.5cm}|P{1.9cm}|}
\hline
\cab{Atividade} & \cab{Riscos potenciais} & \cab{Medidas de mitigação} &
\cab{Instrumento} & \cab{Acompanha\-mento} \\ \hline
Planejamento e gestão &
Exclusão de grupos vulneráveis ou comunidades tradicionais nos diagnósticos rurais &
Consultas públicas e metodologias participativas para seleção das áreas &
PEPI & IAT \\ \hline
Sistemas de abastecimento e distribuição &
Contaminação de aquíferos; impactos temporários de obra (ruído, poeira e erosão local) &
Adoção de normas técnicas de perfuração &
MGAS & IAT \\ \hline
Operação e manutenção &
Falta de sustentabilidade técnica; conflitos locais pelo uso da água &
Treinamento das associações locais; definição de modelos de gestão e tarifa social &
MGAS e Plano de Gestão Social & IAT e prefeituras \\ \hline
Sistemas de tratamento de esgoto &
Manejo inadequado de lodo; contaminação do solo por instalação incorreta &
Capacitação das Unidades de Produção Familiar (UPF); observância da RDC 49/2013 &
MGAS e protocolos de assistência técnica & IAT / IDR-Paraná \\ \hline
Relações de trabalho nas obras &
Acidentes de trabalho ou condições inadequadas para operários e técnicos &
Protocolos de saúde e segurança do trabalho &
PGR & IAT \\ \hline
\end{tabularx}
\fonte{Adaptado de \citeonline[Quadro 26]{mop2026}.}
\end{quadro}

\pendencia{Conferir a citação da RDC 49/2013 (Anvisa). Essa resolução trata
da regularização de atividades de interesse sanitário do microempreendedor
individual e do empreendimento familiar rural, e não do manejo de lodo de
sistemas de esgotamento. Verificar qual norma se pretende citar.}
